% Demand Estimation, Nevo Style: Random-Coefficients Logit on Dominick's Cheese
% LaTeX version of demand_estimation_nevo_style.md. Compiles with pdflatex (or Overleaf).
\documentclass[11pt]{article}

\usepackage[margin=1in]{geometry}
\usepackage[T1]{fontenc}
\usepackage[utf8]{inputenc}
\usepackage{lmodern}
\usepackage{amsmath,amssymb}
\usepackage{booktabs}
\usepackage{tabularx}
\usepackage{array}
\usepackage{enumitem}
\usepackage{xcolor}
\usepackage{listings}
\usepackage[authoryear,round]{natbib}
\usepackage[hidelinks]{hyperref}

\setlist{itemsep=2pt, topsep=4pt}
\newcolumntype{L}{>{\raggedright\arraybackslash}X}
\renewcommand{\arraystretch}{1.15}

\lstset{
  basicstyle=\ttfamily\small,
  columns=fullflexible,
  keepspaces=true,
  breaklines=true,
  frame=single,
  rulecolor=\color{gray!50},
  xleftmargin=1em,
  framexleftmargin=0.5em
}

% Shorthands
\newcommand{\code}[1]{\texttt{#1}}
\newcommand{\family}{\text{family}}
\newcommand{\sigconst}{\sigma_{\text{constant}}}
\newcommand{\sigprice}{\sigma_{\text{price}}}

\title{Demand Estimation, Nevo Style:\\Random-Coefficients Logit on Dominick's Cheese}
\author{}
\date{}

\begin{document}
\maketitle

This note estimates a random-coefficients logit (BLP) demand system for packaged cheese at
Dominick's Finer Foods. The model follows \citet{nevo2000}. Product definition, quantity units
and market size follow \citet{mw2017}.

\begin{itemize}
  \item Data construction: \code{code/01\_upc.py} (UPC attributes), \code{code/02\_clean.py}
        (store-week panel), \code{code/groups.py} (brand and type groups).
  \item Estimation: \code{code/03\_blp\_nevo.py}, pyBLP 1.3.0.
  \item Diagnostics behind the numbers quoted below: \code{code/diagnostics/}.
  \item Cluster scripts: \code{code/longleaf/}.
  \item Environment: \code{requirements.txt} (Python 3.12, used for the benchmark),
        \code{requirements-py39.txt} (Python 3.9); setup steps in \code{README.md}.
\end{itemize}

%==========================================================================================
\section{Benchmark specification at a glance}\label{sec:glance}

\begin{table}[htbp]
\centering
\small
\caption{Benchmark specification}\label{tab:glance}
\begin{tabularx}{\textwidth}{@{}lL@{}}
\toprule
Component & Benchmark choice \\
\midrule
Market & store $\times$ week: 27,230 markets, 83 stores with demographics, 380 weeks between
         September 1989 and April 1997 \\
Product & brand group (8) $\times$ type group (9) $\times$ size class (5); top 41 products,
          covering 90.5\% of all cheese units \\
Quantity & 8-oz equivalent units (total ounces $\div$ 8) \\
Price & dollars per ounce, volume-weighted over the UPCs in a product \\
Market size & $M_s = 1.5 \times \max_w Q_{sw}$: 1.5 times the store's largest weekly inside
              sales, in 8-oz units (Miller--Weinberg) \\
Outside good & no purchase, plus cheese outside the 41 products \\
Model & random-coefficients logit, \textbf{no nest} \\
Fixed effects & store and week, absorbed \\
Linear characteristics $X_1$ & price + 19 dummies (size class, brand group, type group) \\
Random coefficients $X_2$ & constant, price, family-size package \\
Demographic interactions $\Pi$ & constant $\times$ income, price $\times$ income, family size
                                 $\times$ household size \\
Instruments & 51 candidates: wholesale cost (3), BLP sums (28), demographics $\times$
              characteristics (20); screened to a well-conditioned set (49 in the full sample) \\
Estimator & two-step GMM, nested fixed point (Nevo 2000 appendix) \\
\bottomrule
\end{tabularx}
\end{table}

Command: \code{python code/03\_blp\_nevo.py --every 1} (all other options are the defaults).

%==========================================================================================
\section{Mapping to Nevo (2000) and Miller--Weinberg (2017)}\label{sec:mapping}

\begin{table}[htbp]
\centering
\small
\caption{Comparison with Nevo (2000) and Miller--Weinberg (2017)}\label{tab:mapping}
\begin{tabularx}{\textwidth}{@{}lLLL@{}}
\toprule
 & Nevo (2000): cereal & Miller--Weinberg (2017): beer & This note: cheese \\
\midrule
Market & city $\times$ quarter & region $\times$ month & store $\times$ week \\
Products & 24 brands & $\sim$39 brand $\times$ package size & 41 brand group $\times$ type group
           $\times$ size class \\
Quantity & servings & 144-oz equivalents & 8-oz equivalents \\
Market size & population $\times$ 1 serving/day & 1.5 $\times$ max regional sales & 1.5
              $\times$ max weekly store sales \\
Model & RC logit & RC nested logit & RC logit (nest not identified here,
        Section~\ref{sec:history}) \\
Fixed effects & brand & product + time & store + week \\
Demographics & income, income$^2$, age, child & income & income, household size \\
Price instruments & Hausman & distance $\times$ diesel, post-merger dummy & wholesale cost \\
$\Pi$ instruments & --- & mean income $\times$ characteristics & store demographics $\times$
                    characteristics \\
\bottomrule
\end{tabularx}
\end{table}

%==========================================================================================
\section{Data}\label{sec:data}

\subsection{Products}\label{sec:products}

A product is a cell of three coarse characteristics. Every UPC with a known weight is mapped to
one cell; UPCs in the same cell are merged.

\begin{table}[htbp]
\centering
\small
\caption{Brand groups (8)}\label{tab:brands}
\begin{tabularx}{0.85\textwidth}{@{}lL@{}}
\toprule
Brand group & Brands \\
\midrule
private label & Dominick's, Heritage House, Generic \\
Kraft & Kraft, Cracker Barrel, Kraft Healthy Favorites, Light n' Lively \\
Philadelphia & Philadelphia \\
Velveeta & Velveeta \\
Sargento & Sargento, MooTown \\
County Line & County Line \\
Borden & Borden \\
other national & all other brands, incl.\ Healthy Choice, Land O Lakes, Laughing Cow and
                 brands coded ``Other'' \\
\bottomrule
\end{tabularx}
\end{table}

\paragraph{Type groups (9).} Cream cheese; processed; shredded; cheddar block/sliced;
mozzarella block/sliced; other natural block/sliced; grated parmesan; snack (string, crackers);
specialty.

\paragraph{Size classes (5).} Small ($\le$6 oz); medium (7--10 oz); large (12--16 oz); bulk
($>$16 oz); random weight (store-cut packages labelled ``8 TO 1'', i.e.\ 8 oz--1 lb).

\paragraph{Selection.} Cells are ranked by total units over the sample, and the smallest set
covering 90\% of all cheese units is kept: \textbf{41 products covering 90.5\%}. Count packs
(``5 CT'', no weight; 0.8\% of units) and non-cheese items (0.2\%) are never inside goods. The
full list, with member UPCs and descriptions, is in \code{data/products\_coarse.csv}.

\paragraph{Why coarse products.} The cardinality of every characteristic is at most 10, as
advised. Compared with the earlier 150 brand $\times$ type $\times$ exact-size products (93.7\%
coverage), the panel shrinks from 3.28 million to 1.12 million product-market rows, and
estimation is about three times faster.

\subsection{Quantities, prices and shares}\label{sec:shares}

\begin{itemize}
  \item \textbf{Quantity} $q_{jt}$: total ounces sold $\div$ 8. Random-weight packages are
        assumed to weigh 12 oz, the midpoint of 8 oz--1 lb.
  \item \textbf{Price} $p_{jt}$: total revenue $\div$ total ounces, in \$/oz. For random-weight
        items the shelf price is per pound, so \$/oz = shelf price $\div$ 16.
  \item \textbf{Share} $s_{jt} = q_{jt} / M_s$.
\end{itemize}

The share unit (8-oz equivalents) matches the price unit (\$/oz), as in Miller--Weinberg. Each
choice is ``buy one 8-oz unit of cheese, or not''. Package size enters only as a product
characteristic.

\subsection{Market size}\label{sec:marketsize}

\begin{equation}
  M_s = 1.5 \times \max_w Q_{sw}, \qquad Q_{sw} = \sum_{j} q_{jsw}.
\end{equation}

$M_s$ is computed over all weeks of the full panel before any subsampling, so subsample and
full-sample runs share the same market size. Across stores, $M_s$ ranges from 6,373 to 24,998
8-oz units per week (mean 14,282). Inside shares in the full sample: mean 0.283, median 0.271,
5th percentile 0.168, 95th percentile 0.431, maximum 0.667 ($=1/1.5$ by construction, reached
in each store's peak week).

\paragraph{Caveat.} In 62 of 93 stores the maximum week is week 223 (December 16, 1993), a
chain-wide pre-Christmas promotion: chain sales were 2.6 times a normal week, and Philadelphia
8-oz cream cheese sold 15 times its usual volume. \code{--mw-quantile 0.99} replaces the
maximum by the 99th percentile as a robustness check.

An alternative market size, weekly checkout transactions from \code{ccount.dta}, is available
with \code{--market-size traffic}.

\subsection{Other variables in the panel}\label{sec:othervars}

\begin{table}[htbp]
\centering
\small
\caption{Other panel variables}\label{tab:othervars}
\begin{tabularx}{\textwidth}{@{}lL@{}}
\toprule
Variable & Definition \\
\midrule
\code{wcost} & wholesale cost \$/oz $= p \times (1 - \text{PROFIT}/100)$, volume-weighted;
               margins outside $(0, 100)$ treated as missing (0.05\% of rows) \\
\code{promo\_B}, \code{promo\_S}, \code{promo\_BS} & 1 if any UPC in the product carried a
               bonus-buy (B) or simple price reduction (S) code \\
\code{promo\_BS\_share} & share of the product's units sold under B or S \\
\code{promo\_B\_share}, \code{promo\_S\_share} & share of the product's units sold under B,
               under S \\
\code{promo\_B\_fw}, \code{promo\_S\_fw} & fixed-weight share of the product's UPCs on B, on S
               (Section~\ref{sec:promo-measures}) \\
\code{promo\_rate\_oth\_*} & share of the \emph{other} products on B/S in the same market,
               type group, brand group or size class \\
\bottomrule
\end{tabularx}
\end{table}

\subsection{Descriptive statistics (estimation sample)}\label{sec:desc}

\paragraph{Sample.} 27,230 store-weeks (83 stores $\times$ 380 weeks, September 14, 1989 to
April 30, 1997); 41 products; 1,087,156 product-market observations; 39.9 products per market
on average (minimum 28).

\begin{table}[htbp]
\centering
\small
\caption{Product-market variables}\label{tab:desc}
\begin{tabular}{@{}lrrrrr@{}}
\toprule
Variable & Mean & SD & 10th pct. & Median & 90th pct. \\
\midrule
price (\$/oz) & 0.245 & 0.098 & 0.162 & 0.223 & 0.373 \\
wholesale cost (\$/oz) & 0.161 & 0.070 & 0.101 & 0.147 & 0.238 \\
retail margin, $1 - \text{wcost}/p$ & 0.340 & 0.101 & 0.229 & 0.339 & 0.464 \\
on B/S promotion (any UPC) & 0.318 & 0.466 & 0 & 0 & 1 \\
share of units sold on B/S & 0.190 & 0.331 & 0 & 0 & 0.840 \\
product share $s_{jt}$ & 0.0071 & & & 0.0034 & \\
inside share $\sum_j s_{jt}$ (per market) & 0.283 & & 0.168 & 0.271 & \\
\bottomrule
\end{tabular}
\end{table}

Promotion frequency is higher here than at the UPC level (19\%) because a product is flagged
when \emph{any} of its UPCs is on promotion.

\paragraph{Store demographics (83 stores).} Median household income \$41,578 (log income: mean
10.62, SD 0.28 across stores); average household size: mean 2.66, SD 0.26.

\begin{table}[htbp]
\centering
\small
\caption{Volume shares, median price and promotion frequency by characteristic}
\label{tab:bychar}
\begin{tabular}{@{}lrrrr@{}}
\toprule
 & Products & Volume share & Median \$/oz & Share of weeks on B/S \\
\midrule
\multicolumn{5}{@{}l}{\emph{Brand group}} \\
private label & 13 & 39.1\% & 0.195 & 0.27 \\
Kraft & 12 & 27.5\% & 0.257 & 0.34 \\
Philadelphia & 3 & 15.3\% & 0.173 & 0.28 \\
Sargento & 4 & 6.0\% & 0.304 & 0.29 \\
Velveeta & 2 & 5.9\% & 0.177 & 0.43 \\
other national & 4 & 3.5\% & 0.289 & 0.40 \\
County Line & 2 & 1.5\% & 0.266 & 0.33 \\
Borden & 1 & 1.2\% & 0.202 & 0.47 \\
\addlinespace
\multicolumn{5}{@{}l}{\emph{Type group}} \\
processed & 9 & 34.2\% & 0.192 & 0.38 \\
cream cheese & 5 & 21.0\% & 0.170 & 0.32 \\
shredded & 9 & 17.1\% & 0.257 & 0.28 \\
cheddar block/sliced & 4 & 10.2\% & 0.231 & 0.34 \\
other natural block/sliced & 4 & 9.7\% & 0.235 & 0.34 \\
mozzarella block/sliced & 3 & 3.8\% & 0.198 & 0.19 \\
grated parmesan & 3 & 2.8\% & 0.408 & 0.32 \\
snack (string, crackers) & 3 & 0.8\% & 0.354 & 0.24 \\
specialty & 1 & 0.4\% & 0.422 & 0.52 \\
\addlinespace
\multicolumn{5}{@{}l}{\emph{Size class}} \\
medium (7--10 oz) & 17 & 41.0\% & 0.251 & 0.41 \\
large (12--16 oz) & 10 & 38.5\% & 0.197 & 0.42 \\
random weight & 4 & 14.0\% & 0.193 & 0.00 \\
small ($\le$6 oz) & 9 & 3.5\% & 0.332 & 0.17 \\
bulk ($>$16 oz) & 1 & 3.0\% & 0.162 & 0.32 \\
\bottomrule
\end{tabular}
\end{table}

Volume shares are of the 41 inside products' total 8-oz volume. Price per ounce falls with
package size (small \$0.33, medium \$0.25, large \$0.20, bulk \$0.16), and random-weight
store-cut cheese is almost never on B/S promotion.

%==========================================================================================
\section{Model}\label{sec:model}

\subsection{Utility}\label{sec:utility}

\begin{equation}\label{eq:utility}
u_{ijt} = \underbrace{\alpha p_{jt} + x_j\beta + \lambda_s + \tau_w + \xi_{jt}}_{\delta_{jt}}
+ \underbrace{[\,1,\ p_{jt},\ \family_j\,]\,(\Pi D_i + \Sigma\nu_i)}_{\mu_{ijt}}
+ \varepsilon_{ijt},
\end{equation}
with the outside good normalized to $u_{i0t} = \varepsilon_{i0t}$ and $\varepsilon$ i.i.d.\
type-I extreme value.

\begin{itemize}
  \item $\delta_{jt}$ is the mean utility: price, the 19 characteristic dummies $x_j$, store
        fixed effects $\lambda_s$, week fixed effects $\tau_w$ and the structural error
        $\xi_{jt}$.
  \item $\mu_{ijt}$ is the consumer-specific deviation. $[1, p_{jt}, \family_j]$ are the three
        characteristics with random coefficients: the constant (taste for buying cheese at
        all), price (price sensitivity) and the family-size dummy (large or bulk package).
\end{itemize}

\subsection{Linear part $X_1$}\label{sec:x1}

Price, plus 19 dummies: 4 size classes, 7 brand groups and 8 type groups, each relative to a
base category. Store and week fixed effects are absorbed.

\subsection{Random coefficients}\label{sec:rc}

$\Sigma$ is diagonal with three entries (constant, price, family size). $\Pi$ keeps three
interactions (Table~\ref{tab:pi}).

\begin{table}[htbp]
\centering
\small
\caption{Nonzero entries of $\Pi$}\label{tab:pi}
\begin{tabular}{@{}lcc@{}}
\toprule
$X_2$ \textbackslash{} $D$ & income & household size \\
\midrule
constant & \checkmark & 0$^\dagger$ \\
price & \checkmark & --- \\
family size & --- & \checkmark \\
\bottomrule
\multicolumn{3}{@{}p{0.6\textwidth}@{}}{\footnotesize $^\dagger$ Household size is constant
within a store, so constant $\times$ household size would be a store fixed effect.}
\end{tabular}
\end{table}

Each entry has a direct interpretation: price $\times$ income measures how price sensitivity
changes with income, the core interaction in both \citet{nevo2000} and \citet{mw2017};
constant $\times$ income measures how the propensity to buy cheese changes with income; family
size $\times$ household size measures whether larger households prefer larger packages.

\paragraph{Why only three.} Store demographics are highly collinear across the 83 stores:
$\text{corr}(\text{household size}, \text{share of children}) = 0.80$ and
$\text{corr}(\text{share of children}, \text{share of elderly}) = -0.71$. With ten
interactions, the estimates were large, offset each other in sign, and were all insignificant
(Section~\ref{sec:history}). Miller--Weinberg use income only.

\code{--pi full} restores ten interactions: constant $\times$ (income, child, elderly), price
$\times$ (income, income$^2$, household size, child, elderly), family size $\times$ (household
size, child).

\subsection{Fixed effects and random coefficients}\label{sec:fe}

Store fixed effects and random coefficients coexist. Fixed effects shift the \emph{level} of
mean utility; random coefficients govern \emph{substitution}, through price and product
characteristics, which vary across products within a store. Only an interaction of the
constant with a demographic that is constant within a store (here household size) is absorbed
by store fixed effects. \citet{nevo2000} likewise combines brand fixed effects with random
coefficients.

\subsection{Simulated consumers}\label{sec:agents}

Each store has 20 simulated consumers, reused in every week (Table~\ref{tab:agents}).

\begin{table}[htbp]
\centering
\small
\caption{Simulated consumer draws}\label{tab:agents}
\begin{tabularx}{\textwidth}{@{}lL@{}}
\toprule
Variable & Construction \\
\midrule
\code{income} & log income $\sim N(\text{income}_s, s_s^2)$; $\text{income}_s$ is the store's
                log median household income, and $s_s$ matches the dollar standard deviation
                \code{incsigma} under a lognormal ($s_s \approx 0.5$) \\
\code{hhsize} & store average household size, the same for every consumer in the store \\
\code{nodes0}--\code{nodes2} & $\nu_i \sim N(0, I_3)$ \\
\bottomrule
\end{tabularx}
\end{table}

Demographics are demeaned over all agents, so $(\alpha, \beta)$ are the tastes of the average
consumer. (\code{income\_sq}, \code{child}, \code{elderly} are also drawn; they are used only
with \code{--pi full}.)

%==========================================================================================
\section{Instruments}\label{sec:iv}

There are 51 candidate excluded instruments (Table~\ref{tab:iv}).

\begin{table}[htbp]
\centering
\small
\caption{Candidate excluded instruments}\label{tab:iv}
\begin{tabularx}{\textwidth}{@{}lcLl@{}}
\toprule
Group & Count & Construction & Mainly identifies \\
\midrule
Wholesale cost & 3 & own wholesale cost \$/oz; mean wholesale cost of rival products in the
                     market; mean wholesale cost of rival products in the same type group
                   & price \\
BLP sums & 28 & sums of size-class (5) and type-group (9) dummies over the firm's other
                products and over rival firms' products & random coefficients \\
Demographic & 20 & store mean income, household size, share of children, share of elderly
                   (standardized) $\times$ 4 size-class dummies and $\times$ the number of
                   rival products & $\Pi$ \\
\bottomrule
\end{tabularx}
\end{table}

\paragraph{Screening.} pyBLP builds the 2SLS weighting matrix from
$Z = [X_1, \text{excluded instruments}]$ after absorbing the fixed effects, so redundancy is
checked in that space: instruments are within-transformed by store and week, the 19 dummies are
partialled out, and columns are kept from a pivoted QR with relative tolerance $10^{-4}$, then
standardized. The number kept depends on the sample: 49 in the full sample, 46 on the quarterly
subsample (2,170 markets), 35 on the 175-market smoke test. The condition number of the
standardized $[X_1, Z]$ is 763 in the full sample.

\subsection{Wholesale cost}\label{sec:wcost}

Dominick's records each UPC's gross margin \code{PROFIT}, so wholesale cost is
$p_{jt}(1 - \text{PROFIT}_{jt}/100)$. It is a supply-side cost shifter that varies by product
and week, the analogue of Miller--Weinberg's distance $\times$ diesel price:
\begin{itemize}
  \item within a UPC-store, its time-series variation is as large as that of the retail price
        (standard deviation of logs 0.099 vs 0.100), with a correlation of only 0.36;
  \item it is identical across stores in a given UPC-week (a chain-wide cost), so it survives
        week fixed effects through cross-product variation;
  \item a 1-cent increase in wholesale cost raises the retail price by 1.27 cents on average.
\end{itemize}

\paragraph{Validity concerns.} \code{PROFIT} is based on Dominick's average acquisition cost,
not the current replacement cost. Manufacturer trade deals can lower the wholesale cost and
trigger a retail promotion at the same time, and promotions also raise demand directly through
displays and feature advertising. Both would correlate the instrument with $\xi_{jt}$. It is
still the best cost shifter available in these data.

\subsection{First-stage strength}\label{sec:firststage}

Table~\ref{tab:firststage} reports the partial $R^2$ and $F$ of the excluded instruments for
price, after partialling out store and week fixed effects and the 19 dummies (quarterly
subsample, 2,170 markets).

\begin{table}[htbp]
\centering
\small
\caption{First-stage strength for price}\label{tab:firststage}
\begin{tabular}{@{}lll@{}}
\toprule
Instruments & 150-product panel & 41-product panel \\
\midrule
BLP sums only & $R^2 = 0.037$, $F = 377$ & $R^2 = 0.006$, $F = 23$ \\
demographic only & $R^2 = 0.0005$, $F = 6.6$ & $R^2 = 0.001$, $F = 4.6$ \\
own wholesale cost & --- & $R^2 = \mathbf{0.835}$, $F = 438{,}441$ \\
all 46 & --- & $R^2 = 0.839$, $F = 9{,}775$ \\
\bottomrule
\end{tabular}
\end{table}

\paragraph{Why BLP instruments alone fail on the 41-product panel.} With time-invariant
characteristics and store, week and characteristic fixed effects, BLP sums vary only when the
choice set changes. The 41 products sell in almost every store-week (standard deviation of the
number of products per market: 0.7, vs 4.4 with 150 products), so the BLP sums are nearly
constant within a product and carry almost no information about price. A time-varying cost
shifter is required, which is what Miller--Weinberg use.

\subsection{Instruments not used}\label{sec:notused}

\begin{itemize}
  \item \textbf{Hausman} (prices of the same product in other stores): all stores are in
        Chicago and prices are set chain-wide, so other-store prices share the demand shocks.
  \item \textbf{Own promotion flags (B/S):} conditional on price, promotion weeks sell about
        26\% more, so promotions shift demand directly and fail the exclusion restriction.
\end{itemize}

%==========================================================================================
\section{Specification history and diagnostics}\label{sec:history}

All smoke tests use weeks 100, 200 and 300: 175 markets, 83 stores.

\begin{table}[htbp]
\centering
\footnotesize
\setlength{\tabcolsep}{4pt}
\caption{Specification history (smoke tests, 175 markets)}\label{tab:history}
\begin{tabularx}{\textwidth}{@{}clllllrrL@{}}
\toprule
Run & Prod. & Shares, $M$ & Nest & Instruments & Price coef. & Med.\ own el. & Pos.\ el. &
Notes \\
\midrule
1 & 150 & packages, transactions & no & BLP & $-11.3$ (7.2) & $-5.32$ & 10.9\% &
    10 $\Pi$, large offsetting signs \\
2 & 41 & 8-oz, MW & yes & BLP + demog. & $+0.50$ (1.2) & $+11.0$ & 100\% &
    $\rho = 0.99$ (bound); negative step-1 objective \\
3 & 41 & 8-oz, MW & yes & + wholesale cost & $-0.06$ (0.6) & $-2.17$ & 17.3\% &
    $\rho = 0.99$ (bound) \\
4 & 41 & 8-oz, MW & no & + wholesale cost & $-12.9$ (2.4) & $-2.83$ & 0\% &
    10 $\Pi$, all insignificant \\
5 & 41 & 8-oz, MW & no & + wholesale cost & $-19.4$ (4.0) & $-3.40$ & 1.0\% &
    same as 4, different collinear column kept; unstable \\
6 & 41 & 8-oz, MW & no & + wholesale cost & $-10.7$ (5.5) & $-2.40$ & 0\% & 3 $\Pi$ \\
\textbf{7} & \textbf{41} & \textbf{8-oz, MW} & \textbf{no} & \textbf{+ w.\ cost, screened} &
    $\mathbf{-10.7}$ \textbf{(5.0)} & $\mathbf{-2.38}$ & \textbf{0\%} &
    \textbf{benchmark: 3 $\Pi$; weighting-matrix condition number $2\times10^6$} \\
\bottomrule
\end{tabularx}
\end{table}

\paragraph{Lessons.}
\begin{enumerate}
  \item \textbf{Wholesale cost is essential} once products are coarse (runs 2 $\to$ 3).
  \item \textbf{The nest is not identified here} (runs 3 vs 4). The nesting parameter needs
        variation that moves total inside share separately from shares among inside goods,
        e.g.\ changes in the number of products. The choice set is nearly fixed and $M_s$ is
        fixed per store, so $\rho$ runs to its upper bound of 0.99. Miller--Weinberg also
        report a weak first stage for their nesting parameter.
  \item \textbf{Ten $\Pi$ interactions are not identified} from 83 stores with collinear
        demographics (runs 4--5); three interactions are (run 6).
\end{enumerate}

\paragraph{Numerical issue, resolved in run 7.} In runs 1--6 the condition number of the GMM
weighting matrix was of order $10^{18}$--$10^{19}$. The instruments had been screened for
collinearity \emph{before} absorbing the fixed effects and without partialling out the dummies,
so near-dependent columns survived. On Longleaf (newer numpy/scipy) one more such column
survived, and pyBLP failed with ``Failed to compute the 2SLS weighting matrix''. Screening in
the absorbed, dummy-partialled space (Section~\ref{sec:iv}) brings the condition number to
about $2\times10^6$, with essentially unchanged estimates.

%==========================================================================================
\section{Benchmark results}\label{sec:results}

Both runs use the defaults in Section~\ref{sec:glance}, on Longleaf, one core each
(Table~\ref{tab:runs}).

\begin{table}[htbp]
\centering
\small
\caption{Benchmark runs}\label{tab:runs}
\begin{tabular}{@{}lll@{}}
\toprule
 & Quarterly subsample & Full sample \\
\midrule
Markets / stores / weeks & 2,170 / 83 / 30 & 27,230 / 83 / 380 \\
Product-market observations & 86,595 & 1,087,156 \\
Excluded instruments after screening & 46 & 49 \\
Run time / peak memory & 5.5 min / 3.1 GB & 64 min / 37 GB \\
\bottomrule
\end{tabular}
\end{table}

Both GMM steps converge. In the full sample, step 2 stops after 18 iterations, when the
objective improves by less than $10^{-5}$ (relative), with a projected gradient norm of 0.11
on an objective of $1.6\times10^4$. The reduced Hessian has one zero eigenvalue, from
$\sigconst$ sitting at its lower bound of 0. The weighting-matrix condition number is
$2.6\times10^6$.

\subsection{Parameter estimates}\label{sec:estimates}

\begin{table}[htbp]
\centering
\small
\caption{Nonlinear parameters and mean price coefficient}\label{tab:estimates}
\begin{tabular}{@{}lccc@{}}
\toprule
Parameter & Smoke test (175) & Quarterly (2,170) & \textbf{Full (27,230)} \\
\midrule
price (mean $\alpha$) & $-10.66$ (4.97) & $-7.46$ (1.52) & $\mathbf{-8.21}$ \textbf{(0.45)} \\
$\sigma$ constant & 0.00 (34.1) & 0.00 (0.74) & \textbf{0.00 (0.21)} \\
$\sigma$ price & 0.00 (9.39) & 2.65 (0.87) & \textbf{2.93 (0.24)} \\
$\sigma$ family size & 0.00 (1.65) & 0.32 (0.34) & \textbf{0.33 (0.10)} \\
$\pi$ constant $\times$ income & 0.21 (53.6) & 0.56 (0.83) & \textbf{0.31 (0.25)} \\
$\pi$ price $\times$ income & $-0.46$ (3.82) & $-0.20$ (0.27) &
    $\mathbf{-0.12}$ \textbf{(0.08)} \\
$\pi$ family size $\times$ household size & 0.33 (0.10) & 0.38 (0.03) &
    \textbf{0.36 (0.01)} \\
\bottomrule
\multicolumn{4}{@{}l}{\footnotesize Number of markets in parentheses in the header. Standard
errors in parentheses. Price is in dollars per ounce.}
\end{tabular}
\end{table}

\begin{table}[htbp]
\centering
\small
\caption{Linear characteristics, full sample. Base categories: bulk size, Borden, cheddar
block/sliced}\label{tab:linear}
\begin{tabular}{@{}lr@{\hspace{1.5em}}lr@{\hspace{1.5em}}lr@{}}
\toprule
\multicolumn{2}{@{}l}{Size class} & \multicolumn{2}{l}{Brand group} &
\multicolumn{2}{l@{}}{Type group} \\
\midrule
small & $-1.17$ (0.03) & Sargento & 1.68 (0.01) & specialty & 1.74 (0.03) \\
medium & 0.19 (0.03) & Philadelphia & 1.59 (0.01) & grated parmesan & 0.35 (0.03) \\
large & 0.34 (0.01) & private label & 1.19 (0.02) & processed & $-0.04$ (0.01) \\
random weight & 0.26 (0.02) & Kraft & 0.99 (0.01) & other natural block/sliced &
    $-0.05$ (0.00) \\
 & & Velveeta & 0.87 (0.01) & snack & $-0.18$ (0.02) \\
 & & other national & 0.06 (0.01) & shredded & $-0.33$ (0.01) \\
 & & County Line & $-0.08$ (0.01) & cream cheese & $-0.79$ (0.02) \\
 & & & & mozzarella block/sliced & $-1.22$ (0.00) \\
\bottomrule
\end{tabular}
\end{table}

\subsection{Own-price elasticities (full sample)}\label{sec:elasticities}

Across all product-markets: 5th percentile $-2.41$, 25th $-1.80$, \textbf{median $-1.50$},
75th $-1.28$, 95th $-1.00$. No product-market has a positive own-price elasticity.

\begin{table}[htbp]
\centering
\small
\caption{Most and least elastic products (median over markets)}\label{tab:elast}
\begin{tabular}{@{}lr@{\hspace{2em}}lr@{}}
\toprule
Most elastic & & Least elastic & \\
\midrule
Kraft grated parmesan, small & $-2.41$ & private label cream cheese, medium & $-0.92$ \\
other national specialty, small & $-2.28$ & Philadelphia cream cheese, medium & $-1.10$ \\
Kraft grated parmesan, medium & $-2.22$ & private label processed, large & $-1.14$ \\
Sargento snack, small & $-2.14$ & Velveeta processed, bulk & $-1.16$ \\
other national snack, small & $-2.03$ & Philadelphia cream cheese, large & $-1.25$ \\
Kraft snack, small & $-1.97$ & private label processed, random weight & $-1.28$ \\
\bottomrule
\end{tabular}
\end{table}

Full product-level results: \code{output/nevo\_own\_elasticities\_coarse\_every1\_*.csv}.

\subsection{Interpretation}\label{sec:interp}

\begin{enumerate}
  \item \textbf{Stable across samples.} The quarterly and full-sample estimates are close; the
        full sample mainly shrinks the standard errors (price: 1.52 $\to$ 0.45).
  \item \textbf{Price.} The mean price coefficient is $-8.2$. With $\sigprice = 2.93$, the
        price coefficients of the simulated consumers have mean $-8.2$ and standard deviation
        2.9; only 0.26\% of consumers have a positive price coefficient, so the model no
        longer produces positive own elasticities.
  \item \textbf{Elasticities are moderate}, median $-1.5$, below Miller--Weinberg's beer
        estimates ($-4.3$ to $-6.1$). They are highest for small, high-price-per-ounce
        products (grated parmesan, snacks) and lowest for the large staples (cream cheese,
        processed slices). Because quantity is measured in 8-oz units and the products are
        coarse aggregates, these are elasticities of cell-level volume, which are naturally
        lower than UPC-level elasticities.
  \item \textbf{Heterogeneity.} Random coefficients on price and on family size are
        significant. Family size $\times$ household size is strongly positive (0.36,
        $t \approx 45$): stores with larger households buy relatively more large packages.
        Constant $\times$ income is not significant. Price $\times$ income is negative and only
        marginally significant ($t \approx -1.6$), i.e.\ if anything higher-income stores are
        slightly \emph{more} price sensitive, the opposite of the usual sign. Store
        demographics vary little across 83 stores, which limits what $\Pi$ can pick up.
  \item \textbf{$\sigconst = 0$.} No unobserved heterogeneity in the taste for cheese relative
        to the outside good beyond what store fixed effects and income capture.
  \item \textbf{Linear characteristics.} Among brand groups, Sargento, Philadelphia and private
        label have the highest mean utility given price; County Line and the ``other
        national'' group the lowest. Small packages are least preferred; specialty cheese has
        the highest type effect.
\end{enumerate}

\subsection{Smoke test of the benchmark (run 7, 175 markets)}\label{sec:smoke}

Run time 0.2 minutes; both GMM steps converge.

\begin{table}[htbp]
\centering
\small
\caption{Benchmark specification on the smoke-test sample}\label{tab:smoke}
\begin{tabular}{@{}lrr@{}}
\toprule
Parameter & Estimate & SE \\
\midrule
price (mean $\alpha$) & $-10.66$ & 4.97 \\
$\sigma$ constant & 0.000 & 34.1 \\
$\sigma$ price & 0.000 & 9.39 \\
$\sigma$ family size & 0.000 & 1.65 \\
$\pi$ constant $\times$ income & 0.207 & 53.6 \\
$\pi$ price $\times$ income & $-0.459$ & 3.82 \\
$\pi$ family size $\times$ household size & \textbf{0.329} & \textbf{0.097} \\
\bottomrule
\end{tabular}
\end{table}

Own-price elasticities (5th, 25th, 50th, 75th, 95th percentiles): $-4.32$, $-2.99$, $-2.38$,
$-1.99$, $-1.50$. No product-market has a positive own-price elasticity.

Family size $\times$ household size is already positive and significant here
($t \approx 3.4$). All three $\sigma$ sit at their lower bound of 0 in this small sample; with
more markets, $\sigprice$ and $\sigma_{\text{family size}}$ move away from 0 and become
significant (Section~\ref{sec:estimates}).

%==========================================================================================
\section{Promotion controls}\label{sec:promo}

\subsection{Motivation}\label{sec:promo-motivation}

Dominick's flags each UPC-store-week with a deal code. Two codes matter for cheese:
\textbf{B} (bonus buy) and \textbf{S} (simple price reduction). The question here is whether
consumers respond to the \emph{information} that a product is on promotion (shelf tags,
signage, display), over and above the price cut itself. The price cut is already captured by
$p_{jt}$, so a promotion variable in $X_1$ measures the extra utility from the promotion signal
at a given price.

Promotions are observed at the store-week level, not for each shopper, so they enter as
product-market characteristics in the mean utility $\delta_{jt}$. They are treated as
exogenous and therefore also serve as their own instruments.

Note that the Dominick's code B means ``bonus buy'', an in-store featured deal; it is not
necessarily a buy-one-get-one offer.

\subsection{Measures}\label{sec:promo-measures}

A product merges several UPCs, so a promotion can cover part of a product.
Table~\ref{tab:promo-measures} lists three measures, each with one column for B and one for S.

\begin{table}[htbp]
\centering
\small
\caption{Promotion measures}\label{tab:promo-measures}
\begin{tabularx}{\textwidth}{@{}lLL@{}}
\toprule
Option & Definition & Interpretation \\
\midrule
\code{--promo any} & 1 if any UPC of the product carries the code & a promotion tag is visible
                     for the product \\
\code{--promo fw} & $\sum_u w_{us}\,1[u \text{ on deal}] / \sum_u w_{us}$ over the product's
                    UPCs sold that week; $w_{us}$ is UPC $u$'s share of the product's units in
                    store $s$ over the whole sample & probability that a typical shopper of the
                    product meets a promoted UPC \\
\code{--promo fwdummy} & $1[\text{fw} \ge 0.5]$ (\code{--promo-cut}) & most of the product is
                         on promotion \\
\bottomrule
\end{tabularx}
\end{table}

\paragraph{Why fixed weights.} The obvious measure, the share of this week's units sold on the
deal (\code{promo\_B\_share}, \code{promo\_S\_share}, available as \code{--promo share}), has
this week's quantity in the denominator. A promotion raises the promoted UPC's sales, so that
share moves mechanically with the demand shock $\xi_{jt}$. The fixed weights $w_{us}$ do not
change from week to week, so \code{fw} reflects only \emph{which} UPCs are promoted. The two
measures are highly correlated (0.97 for B, 0.96 for S), but only \code{fw} can be treated as
an exogenous characteristic.

\begin{table}[htbp]
\centering
\small
\caption{Promotion variables, estimation sample (1,087,156 product-markets)}
\label{tab:promo-desc}
\begin{tabular}{@{}lcc@{}}
\toprule
 & B & S \\
\midrule
product-markets with any UPC on the deal & 28.0\% & 5.0\% \\
mean of \code{fw} & 0.148 & 0.028 \\
mean of \code{fw} when positive & 0.53 & 0.57 \\
share with \code{fw} $\ge 0.5$ & 13.9\% & 2.8\% \\
log price relative to product-store mean, on vs off the deal & $-0.048$ vs $+0.019$ &
    $-0.084$ vs $+0.004$ \\
correlation of \code{fw} with that log price deviation & $-0.30$ & $-0.19$ \\
\bottomrule
\end{tabular}
\end{table}

B and S rarely coincide (1.1\% of product-markets). Random-weight cheese is never on B or S, and
small packages rarely are. Promotions coincide with lower prices, but far from perfectly, so
price and the promotion signal are separately identified.

\subsection{Results}\label{sec:promo-results}

All other settings are the benchmark defaults (Section~\ref{sec:glance}).
Table~\ref{tab:promo-q13} compares the three measures on the quarterly subsample;
Table~\ref{tab:promo-full} reports the preferred measure, \code{fw}, on the full sample.

\begin{table}[htbp]
\centering
\small
\caption{Promotion controls, quarterly subsample (2,170 markets, 46 instruments)}
\label{tab:promo-q13}
\begin{tabular}{@{}lcccc@{}}
\toprule
 & none (benchmark) & any & fw & fwdummy \\
\midrule
price (mean $\alpha$) & $-7.46$ (1.52) & $-4.62$ (1.58) & $-4.54$ (1.63) & $-5.05$ (1.60) \\
promo B & --- & 0.45 (0.02) & 0.47 (0.03) & 0.31 (0.02) \\
promo S & --- & 0.76 (0.03) & 1.10 (0.06) & 0.89 (0.04) \\
$\sigma$ price & 2.65 (0.87) & 1.93 (0.89) & 1.89 (0.92) & 2.05 (0.90) \\
$\sigma$ family size & 0.32 (0.34) & 0.36 (0.44) & 0.32 (0.45) & 0.32 (0.42) \\
$\pi$ price $\times$ income & $-0.20$ (0.27) & $-0.10$ (0.25) & $-0.06$ (0.25) &
    $-0.12$ (0.25) \\
$\pi$ family size $\times$ household size & 0.38 (0.03) & 0.38 (0.03) & 0.38 (0.03) &
    0.38 (0.03) \\
median own elasticity & $-1.40$ & $-0.88$ & $-0.87$ & $-0.96$ \\
own elasticity, 5th--95th pct. & $-2.26$ to $-0.93$ & $-1.46$ to $-0.57$ &
    $-1.43$ to $-0.56$ & $-1.57$ to $-0.62$ \\
\bottomrule
\multicolumn{5}{@{}l}{\footnotesize Standard errors in parentheses.}
\end{tabular}
\end{table}

\begin{table}[htbp]
\centering
\small
\caption{Promotion controls (\code{fw}), full sample (27,230 markets, 49 instruments)}
\label{tab:promo-full}
\begin{tabular}{@{}lcc@{}}
\toprule
 & Benchmark (Section~\ref{sec:results}) & \textbf{With \code{fw} promotion controls} \\
\midrule
price (mean $\alpha$) & $-8.21$ (0.45) & $\mathbf{-6.19}$ \textbf{(0.46)} \\
promo B (fw) & --- & \textbf{0.46 (0.009)} \\
promo S (fw) & --- & \textbf{0.98 (0.016)} \\
$\sigma$ constant & 0.00 (0.21) & 0.00 (0.23) \\
$\sigma$ price & 2.93 (0.24) & 2.50 (0.24) \\
$\sigma$ family size & 0.33 (0.10) & 0.35 (0.11) \\
$\pi$ constant $\times$ income & 0.31 (0.25) & 0.12 (0.25) \\
$\pi$ price $\times$ income & $-0.12$ (0.08) & $-0.08$ (0.07) \\
$\pi$ family size $\times$ household size & 0.36 (0.008) & 0.36 (0.009) \\
\addlinespace
own elasticity, 5th / 25th pct. & $-2.41$ / $-1.80$ & $-1.83$ / $-1.36$ \\
own elasticity, \textbf{median} & $\mathbf{-1.50}$ & $\mathbf{-1.13}$ \\
own elasticity, 75th / 95th pct. & $-1.28$ / $-1.00$ & $-0.96$ / $-0.75$ \\
products with median own elasticity above $-1$ & 1 of 41 & 11 of 41 \\
run time / peak memory (Longleaf, one core) & 64 min / 37 GB & 101 min / 40 GB \\
\bottomrule
\multicolumn{3}{@{}l}{\footnotesize Standard errors in parentheses.}
\end{tabular}
\end{table}

Command: \code{python code/03\_blp\_nevo.py --every 1 --promo fw}. Both GMM steps converge
(30 optimizer iterations in total); the step-2 projected gradient norm is 0.05 and the
weighting-matrix condition number $2.1\times10^6$. As in the benchmark, the reduced Hessian
has one zero eigenvalue, from $\sigconst$ at its bound. The linear characteristics barely move,
except that the type effects of grated parmesan, specialty and snacks (high prices per ounce)
fall by 0.15--0.27 and that of cream cheese rises by 0.15.

\subsection{Interpretation}\label{sec:promo-interp}

\begin{enumerate}
  \item \textbf{The promotion signal matters.} Both coefficients are large and precisely
        estimated in every measure. In dollar terms (coefficient $\div\ |\alpha|$, full
        sample): an S tag on the whole product is worth \$0.16/oz, about 65\% of the mean
        price; a B promotion is worth \$0.07/oz, about 30\%. Simple price reductions carry the
        stronger signal.
  \item \textbf{Price sensitivity was overstated without the controls.} Without them, the
        sales spike in promotion weeks is attributed entirely to the price cut. With them, the
        mean price coefficient falls by a quarter in the full sample ($-8.2 \to -6.2$) and by
        about 40\% in the quarterly subsample; $\sigprice$ falls from 2.9 to 2.5 (0.7\% of
        simulated consumers then have a positive price coefficient).
  \item \textbf{Heterogeneity is unchanged.} Family size $\times$ household size stays at
        0.36; price $\times$ income and constant $\times$ income remain insignificant.
  \item \textbf{The three measures agree.} \code{any} and \code{fw} give almost identical
        results; \code{fwdummy} gives a somewhat smaller B coefficient and a larger price
        coefficient, because it codes partial promotions as zero. \code{fw} is the preferred
        measure: it matches the ``probability that a shopper sees the promotion''
        interpretation and does not move with this week's quantities.
  \item \textbf{Full sample vs quarterly.} The price coefficient is larger in magnitude in the
        full sample ($-6.2$ vs $-4.5$), whereas without promotion controls the two samples
        agreed closely ($-8.2$ vs $-7.5$). With promotion controls, price is identified only
        from cost variation not accompanied by a promotion; a likely reason is that this
        variation is thinner in the quarterly subsample. The full-sample estimates are the
        ones to report.
\end{enumerate}

\paragraph{Caveats.}
\begin{itemize}
  \item \textbf{Elasticities near or below one in magnitude.} The median own elasticity is
        $-1.13$ and 11 of 41 products have a median above $-1$ (cream cheese, processed
        slices, Velveeta). With a 34\% average retail margin, a single-product price setter
        would imply an elasticity around $-3$; inelastic demand is inconsistent with profit
        maximization on that product alone. Multiproduct pricing and loss-leader pricing
        soften this, but the gap suggests the controls may absorb part of the true price
        response: S is itself a price cut, and the volume-weighted product price is a blunter
        measure of the shelf price than the promotion flag is.
  \item \textbf{Instrument interaction.} Manufacturer trade deals lower wholesale cost and
        trigger promotions at the same time. With promotions in $X_1$, that part of the cost
        variation is controlled for, and price is identified from the remaining cost
        variation.
  \item \textbf{Exogeneity.} Promotion timing is a retailer decision and may respond to
        expected demand (e.g.\ holidays; week fixed effects absorb chain-wide timing). If
        promotions are scheduled in high-$\xi$ weeks, their coefficients are biased upward.
\end{itemize}

The benchmark (no controls, median $-1.50$) and this specification (median $-1.13$) bracket
the plausible own-price elasticities; reporting both, with the direction of each bias, is the
honest summary.

%==========================================================================================
\section{Running on Longleaf (UNC)}\label{sec:longleaf}

\begin{enumerate}
  \item Copy code and data from the project root (only \code{code/}, \code{demo.dta} and
        \code{data/*.parquet} are needed):
\begin{lstlisting}[language=bash]
rsync -a --exclude __pycache__ code demo.dta longleaf:empirical_io_ps3/
rsync -a data/panel_week_coarse.parquet longleaf:empirical_io_ps3/data/
\end{lstlisting}
  \item One-time environment setup on a login node: \code{bash code/longleaf/setup.sh}
        (module \code{python/3.12.4}, pyBLP 1.3.0).
  \item Submit:
\begin{lstlisting}[language=bash]
sbatch --partition=general --mem=48g --time=12:00:00 --job-name=BENCH_full \
       --export=ALL,FE=store_week,EVERY=1 code/longleaf/run_nevo.sl
sbatch --partition=general --mem=16g --time=12:00:00 --job-name=BENCH_q13 \
       --export=ALL,FE=store_week,EVERY=13 code/longleaf/run_nevo.sl
\end{lstlisting}
        Extra estimation options are passed through \code{EXTRA}, e.g.\
        \code{--export=ALL,FE=store\_week,EVERY=13,EXTRA="--pi full"}.
  \item Monitor with \code{squeue -u <onyen>} and \code{tail -f logs/<job>\_<id>.out}.
\end{enumerate}

\paragraph{Resources.} Jobs run on one core: \code{pyblp.parallel} was much slower than a
single process on this problem, because there are many small markets and inter-process data
transfer dominates. Actual benchmark runs: quarterly subsample 5.5 minutes and 3.1 GB; full
sample 64 minutes and 37 GB, at about 50 seconds per objective evaluation (step 2: 18 optimizer
iterations, 22 objective evaluations). \code{--mem=48g} on the \code{general} partition is
enough for the full sample.

To fetch results to a laptop (macOS ships an old rsync that accepts one remote path per call):
\begin{lstlisting}[language=bash]
rsync -av --exclude '*.pkl' longleaf:empirical_io_ps3/output/ output/
scp 'longleaf:empirical_io_ps3/logs/BENCH_*' logs/
\end{lstlisting}
Off campus, connect to the UNC VPN first.

%==========================================================================================
\section{Remaining departures and extensions}\label{sec:extensions}

\begin{enumerate}
  \item \textbf{Market size.} $M_s$ is fixed per store and driven by one holiday week in most
        stores. Robustness: \code{--mw-quantile 0.99}, \code{--mw-mult 2}, or
        \code{--market-size traffic}.
  \item \textbf{Nest.} The nesting parameter is not identified; fixing $\rho$ at a few values
        (e.g.\ 0.3, 0.5, 0.7) and comparing elasticities would follow Miller--Weinberg's
        robustness checks.
  \item \textbf{Promotions.} The benchmark has no promotion controls, so the demand effect of
        the promotion signal loads on price. Section~\ref{sec:promo} adds them
        (\code{--promo fw}); a random coefficient on the promotion variables, interacted with
        income, would show which consumers respond to them.
  \item \textbf{Optimal instruments.} \code{results.compute\_optimal\_instruments()} in pyBLP
        \citep{rv2014} may sharpen the estimates of $\Sigma$ and $\Pi$.
  \item \textbf{Stockpiling.} A reduced-form check finds a post-promotion dip of about 8\% in
        the following week (about 15--20\% of the promotion lift), so static elasticities are
        short-run and likely somewhat overstated.
  \item \textbf{Coarse products} mix UPCs within a cell (e.g.\ 12- and 16-oz Kraft Singles,
        regular and reduced-fat), so a promotion on one UPC moves the product's average price
        only partly.
\end{enumerate}

%==========================================================================================
\begin{thebibliography}{9}

\bibitem[Miller and Weinberg(2017)]{mw2017}
Miller, N.~H. and M.~C. Weinberg (2017).
\newblock Understanding the price effects of the MillerCoors joint venture.
\newblock \emph{Econometrica} 85(6), 1763--1791.

\bibitem[Nevo(2000)]{nevo2000}
Nevo, A. (2000).
\newblock A practitioner's guide to estimation of random-coefficients logit models of demand.
\newblock \emph{Journal of Economics \& Management Strategy} 9(4), 513--548.

\bibitem[Reynaert and Verboven(2014)]{rv2014}
Reynaert, M. and F. Verboven (2014).
\newblock Improving the performance of random coefficients demand models: The role of optimal
  instruments.
\newblock \emph{Journal of Econometrics} 179(1), 83--98.

\end{thebibliography}

\end{document}
